\documentclass[11pt]{exam}

\usepackage[interro]{profnsi}
\profnsisetup{
	niveau=Terminale,
	matiere=NSI,
	titre=Interro Python
}

\begin{document}
	\creerDoc
	
    \noindent
    \\
    On souhaite écrire une fonction \texttt{anagramme} qui prend en paramètres deux chaines de caractères \texttt{mot1} et \texttt{mot2} et qui renvoie \texttt{True} si ces deux mots sont des anagrammes, et \texttt{False} sinon.\\
    \textbf{Précision : } Deux mots sont des anagrammes s'ils sont composés exactement des mêmes lettres, avec le même nombre d'occurrences.

    \bigskip
    \textbf{Exemples attendus :}

    \begin{itemize}
        \item \texttt{anagramme("chien", "niche")} renvoie \texttt{True}
        \item \texttt{anagramme("python", "typhon")} renvoie \texttt{True}
        \item \texttt{anagramme("chat", "chats")} renvoie \texttt{False}
    \end{itemize}

    \bigskip

    \lignes{20}

\end{document}
